% ============================================================
% CHƯƠNG 9: Benchmarking — Latency, Scaling, Storage
% ============================================================
\section{Đánh giá hiệu năng}

% ── SLIDE 27: Latency Analysis - Table ──
\begin{frame}{Latency Analysis --- Bảng tổng quan}
\centering
\renewcommand{\arraystretch}{1.1}
{\small
\begin{tabular}{l *{6}{c}}
\toprule
\multirow{2}{*}{\textbf{Phương pháp}} & \multicolumn{2}{c}{\textbf{100K}} & \multicolumn{2}{c}{\textbf{500K}} & \multicolumn{2}{c}{\textbf{1.2M}} \\
\cmidrule(lr){2-3} \cmidrule(lr){4-5} \cmidrule(lr){6-7}
& P50 & P95 & P50 & P95 & P50 & P95 \\
\midrule
Keyword (BM25) & 15 & 56 & 21 & 37 & 71 & 153 \\
Semantic (kNN) & 13 & 34 & 20 & 41 & 191 & 1206 \\
\rowcolor{highlightRow}
\textcolor{accentTeal}{\textbf{Hybrid (RRF)}} & \textcolor{accentTeal}{30} & \textcolor{accentTeal}{53} & \textcolor{accentTeal}{42} & \textcolor{accentTeal}{62} & \textcolor{accentTeal}{\textbf{99}} & \textcolor{accentTeal}{156} \\
\bottomrule
\end{tabular}
}

\vspace{0.05cm}
{\scriptsize (Đơn vị: milliseconds)}

\vspace{0.15cm}

\begin{exampleblock}{Nhận xét về độ trễ}
{\small
\begin{itemize}\itemsep0pt \parskip0pt \parsep0pt
    \item \textbf{Keyword search} duy trì độ trễ rất thấp ở mọi quy mô ($\leq$ 71ms P50).
    \item \textbf{Hybrid Search} đạt hiệu năng ấn tượng với \textbf{99ms P50} ở quy mô 1.2M tài liệu, đáp ứng tốt tiêu chuẩn thời gian thực ($<$ 100ms).
\end{itemize}
}
\end{exampleblock}
\end{frame}

% ── SLIDE 28: Latency Analysis - Charts & Cold Start ──
\begin{frame}{Latency Analysis --- Trực quan hóa \& Thách thức}
\begin{columns}[T]
\begin{column}{0.48\textwidth}
\begin{center}
\includegraphics[width=\textwidth]{latency_comparison.png}

\vspace{0.1cm}
{\scriptsize Độ trễ các phương thức ở quy mô 1.2M}
\end{center}
\end{column}
\begin{column}{0.48\textwidth}

\begin{alertblock}{Thách thức: Semantic Cold Start}
{\small
P95 của Semantic đạt \textbf{1,206ms} ở quy mô 1.2M!\\
\textbf{Nguyên nhân}: HNSW graph (10.4GB) dung lượng lớn phải load từ SSD vào Page Cache trong truy vấn đầu tiên.\\
\textbf{Giải pháp}: Cần warmup cache trước khi phục vụ thực tế.
}
\end{alertblock}

\end{column}
\end{columns}
\end{frame}

% ── SLIDE 29: Latency Scaling ──
\begin{frame}{Latency Scaling --- Tăng trưởng theo quy mô}
\begin{columns}[T]
\begin{column}{0.50\textwidth}
\centering
\includegraphics[width=0.95\textwidth]{latency_scaling.png}

\vspace{0.05cm}
{\scriptsize Độ trễ P50 tăng trưởng theo quy mô dữ liệu}
\end{column}
\begin{column}{0.46\textwidth}

\begin{block}{Phân tích Scaling}
{\scriptsize
\begin{itemize}\itemsep0pt \parskip0pt \parsep0pt
    \item \textbf{Keyword}: Inverted Index $O(1)$, tăng tuyến tính chậm khi scale (15ms $\to$ 71ms).
    \item \textbf{Semantic (HNSW)}: Phức tạp $O(\log N)$, tăng nhanh do đồ thị HNSW lớn \& ảnh hưởng IO.
    \item \textbf{Hybrid}: Bị giới hạn bởi stream chậm hơn, nhưng nhờ chạy song song nên tối ưu chỉ còn 99ms P50.
\end{itemize}
}
\end{block}

\vspace{-0.1cm}

\begin{exampleblock}{Môi trường thử nghiệm}
{\scriptsize
\begin{itemize}\itemsep0pt \parskip0pt \parsep0pt
    \item CPU: Intel Core i7-11800H (8 nhân/16 luồng)
    \item RAM: 16GB DDR4 \quad|\quad SSD: NVMe 512GB
    \item ES: Single node trong Docker trên WSL2 Ubuntu
\end{itemize}
}
\end{exampleblock}

\end{column}
\end{columns}
\end{frame}

% ── SLIDE 30: Storage Analysis - Table ──
\begin{frame}{Storage Analysis --- Dung lượng lưu trữ}
\centering
\renewcommand{\arraystretch}{1.1}
{\small
\begin{tabular}{l c c c}
\toprule
\textbf{Quy mô (Docs)} & \textbf{Số tài liệu} & \textbf{Dung lượng JSONL} & \textbf{Dung lượng ES Index} \\
\midrule
100K & 100,000 & 917 MB & 860 MB \\
500K & 500,000 & 4.6 GB & 4.3 GB \\
1.2M & 1,203,108 & 12 GB & 11.9 GB \\
\bottomrule
\end{tabular}
}

\vspace{0.2cm}

\begin{columns}[T]
\begin{column}{0.48\textwidth}
\begin{block}{Phân tích bộ nhớ}
{\footnotesize
\begin{itemize}\itemsep0pt \parskip0pt \parsep0pt
    \item Tăng trưởng dung lượng \textbf{tuyến tính} theo số lượng tài liệu.
    \item Elasticsearch nén tốt hơn raw JSONL nhờ cấu trúc nén khối của Lucene.
\end{itemize}
}
\end{block}
\end{column}
\begin{column}{0.48\textwidth}
\begin{exampleblock}{Dual-Index ở quy mô 1.2M}
{\footnotesize
\begin{tabular}{l r}
\toprule
\textbf{Index} & \textbf{Dung lượng} \\
\midrule
\texttt{arxiv\_text} (BM25) & 1.5 GB \\
\texttt{arxiv\_vectors} (kNN) & 10.4 GB \\
\textbf{Tổng cộng} & \textbf{11.9 GB} \\
\bottomrule
\end{tabular}
}
\end{exampleblock}
\end{column}
\end{columns}
\end{frame}

% ── SLIDE 31: Storage Analysis - Charts ──
\begin{frame}{Storage Analysis --- Trực quan hóa}
\begin{columns}[T]
\begin{column}{0.50\textwidth}
\begin{center}
\includegraphics[width=\textwidth]{index_size.png}

\vspace{0.1cm}
{\scriptsize Tăng trưởng dung lượng Index theo quy mô}
\end{center}
\end{column}
\begin{column}{0.46\textwidth}
\begin{block}{Đặc điểm dữ liệu Vector}
{\small
\begin{itemize}\itemsep0pt \parskip0pt \parsep0pt
    \item \textbf{87\%} dung lượng lưu trữ của hệ thống là dành cho \textbf{dữ liệu vector} (HNSW graph).
    \item Mỗi vector nhúng: $384 \text{ dims} \times 4 \text{ bytes} = 1{,}536 \text{ bytes}$.
    \item Việc tách thành Dual-Index giúp truy vấn Keyword search không bị tải dung lượng đồ thị HNSW vào bộ nhớ.
\end{itemize}
}
\end{block}
\end{column}
\end{columns}
\end{frame}
